\documentclass[10pt,a4paper]{amsart}
\usepackage[margin=19mm]{geometry}
\usepackage{amsmath,amssymb,mathtools}
\usepackage[scr=rsfs,cal=euler]{mathalfa}
\usepackage{xfrac}
\usepackage[unicode,hidelinks,bookmarks=false]{hyperref}
\newcommand{\ie}{i.e.\ }
\newcommand{\cf}{cf.\ }
\newcommand{\resp}{respectively\ }
\newcommand{\Subsection}{\subsection}
\providecommand{\nobreakdash}{\nobreak\hbox{-}}
\setlength{\emergencystretch}{2em}
\multlinegap0pt
\usepackage{amssymb}
\usepackage[shortlabels]{enumitem}
\usepackage{xcolor}
\newtheorem{lemma}{Lemma}
\title{Operator Michelson Contrast: Logarithmic Parametrisation, Subadditivity, and a Sharp Noncommutativity Threshold}
\author{Irina Nikolaeva, Andrej Novikov}
\date{Source: arXiv:2609.06997v2 (CC BY 4.0)}
\begin{document}
\maketitle

\noindent\emph{Source and licence.} Operator Michelson Contrast: Logarithmic Parametrisation, Subadditivity, and a Sharp Noncommutativity Threshold. Version arXiv:2609.06997v2 (CC BY 4.0), distributed by arXiv under the Creative Commons Attribution 4.0 International licence, http://creativecommons.org/licenses/by/4.0/, as recorded in the arXiv licence metadata for this exact version and shown on the abstract page https://arxiv.org/abs/2609.06997v2. Full licence notice: \texttt{licenses/arxiv-2609.06997-CC-BY-4.0.txt}.\par\smallskip

\noindent\textit{Editorial scope.} Counted unit: the article's lemma lem:SPSQ (source0.tex lines 708--710). The article's complete original proof of it is delivered with it (source0.tex lines 711--753, one \texttt{proof} environment of 43 lines). No mathematics is written, repaired or re-derived: the statement and the proof are the article's own lines, in the article's own notation, and no retained line is substituted or re-flowed.  No further source-local context is needed: every cross-reference of every retained line resolves inside the two delivered spans.  Nothing was omitted from any retained span, so the package carries no omission record.  No retained line contains a citation, so the wrapper carries no bibliography and no bibliography entry.  No retained line contains a figure environment, an image-inclusion command or drawing code, so no image is delivered and the package declares no asset dependency.  Editorial formatting only, in addition: an AMS \texttt{amsart} preamble that restates the article's own declarations and macros used by the retained lines, namely the environments \texttt{lemma} on separate counters, together with the standard packages those lines need.  The retained environments print in this document as Lemma 1 in order of appearance, whereas the article prints the same items with the numbering of its own full document.  The complete native source of the article as distributed in the arXiv e-print for this exact version is bundled unchanged as \texttt{sources/209596/source0.tex}.
\begin{lemma}\label{lem:SPSQ}
For every $\kappa>1$, $S_Q(\kappa) < S_P(\kappa)$.
\end{lemma}

\begin{proof}
First simplify both closed forms. Writing $S_P(\kappa)=\tfrac{2\kappa}{2\kappa+1}[\ln(2\kappa+1)-\ln\kappa]+\tfrac1{2\kappa+1}\ln(2\kappa+1)$ and collecting the $\ln(2\kappa+1)$ terms (their coefficients sum to $\tfrac{2\kappa+1}{2\kappa+1}=1$) gives
\[
S_P(\kappa) = \ln(2\kappa+1) - \frac{2\kappa}{2\kappa+1}\ln\kappa.
\]
An identical manipulation gives
\[
S_Q(\kappa) = \ln(\kappa+2) - \frac{\kappa}{\kappa+2}\ln\kappa.
\]
Define $\varphi(\kappa):=S_P(\kappa)-S_Q(\kappa)$ for $\kappa\ge1$. Then
\[
\varphi(\kappa) = \ln\frac{2\kappa+1}{\kappa+2} - \ln\kappa\left(\frac{2\kappa}{2\kappa+1}-\frac{\kappa}{\kappa+2}\right).
\]
A direct computation gives
\[
\frac{2\kappa}{2\kappa+1}-\frac{\kappa}{\kappa+2} = \kappa\cdot\frac{2(\kappa+2)-(2\kappa+1)}{(2\kappa+1)(\kappa+2)} = \frac{3\kappa}{(2\kappa+1)(\kappa+2)},
\]
so, writing $D(\kappa):=(2\kappa+1)(\kappa+2)=2\kappa^2+5\kappa+2$,
\[
\varphi(\kappa) = \ln\frac{2\kappa+1}{\kappa+2} - \frac{3\kappa\ln\kappa}{D(\kappa)}.
\]
At $\kappa=1$: $\tfrac{2\kappa+1}{\kappa+2}=1$ so the first term vanishes, and $\ln\kappa=0$ kills the second term; hence $\varphi(1)=0$ (consistent with $S_P(1)=S_Q(1)=\ln3$, both computed directly from the simplified forms above).

We show $\varphi'(\kappa)>0$ for $\kappa>1$, which together with $\varphi(1)=0$ gives $\varphi(\kappa)>0$, i.e.\ $S_P(\kappa)>S_Q(\kappa)$, for all $\kappa>1$. Differentiating the first term,
\[
\frac{d}{d\kappa}\ln\frac{2\kappa+1}{\kappa+2} = \frac2{2\kappa+1}-\frac1{\kappa+2} = \frac{2(\kappa+2)-(2\kappa+1)}{D(\kappa)} = \frac{3}{D(\kappa)}.
\]
Differentiating the second term, with $D'(\kappa)=4\kappa+5$,
\[
\frac{d}{d\kappa}\left(\frac{3\kappa\ln\kappa}{D(\kappa)}\right) = \frac{3(\ln\kappa+1)}{D(\kappa)} - \frac{3\kappa\ln\kappa\,D'(\kappa)}{D(\kappa)^2}.
\]
Hence
\[
\varphi'(\kappa) = \frac{3}{D(\kappa)} - \frac{3(\ln\kappa+1)}{D(\kappa)} + \frac{3\kappa\ln\kappa\,D'(\kappa)}{D(\kappa)^2}
= \frac{-3\ln\kappa}{D(\kappa)} + \frac{3\kappa(4\kappa+5)\ln\kappa}{D(\kappa)^2}
= \frac{3\ln\kappa\big[\kappa(4\kappa+5)-D(\kappa)\big]}{D(\kappa)^2}.
\]
Now $\kappa(4\kappa+5)-D(\kappa) = (4\kappa^2+5\kappa) - (2\kappa^2+5\kappa+2) = 2\kappa^2-2 = 2(\kappa-1)(\kappa+1)$, so
\[
\varphi'(\kappa) = \frac{6(\kappa-1)(\kappa+1)\ln\kappa}{D(\kappa)^2}.
\]
For $\kappa>1$, each of $(\kappa-1)$, $(\kappa+1)$, $\ln\kappa$, and $D(\kappa)^2$ is (strictly) positive, so $\varphi'(\kappa)>0$ on $(1,\infty)$. Since $\varphi$ is continuous on $[1,\infty)$ with $\varphi(1)=0$ and $\varphi'>0$ on $(1,\infty)$, we conclude $\varphi(\kappa)>0$, i.e.\ $S_Q(\kappa)<S_P(\kappa)$, for every $\kappa>1$.
\end{proof}

\end{document}
